Con el objetivo de identificar de forma precisa las necesidades del cliente y establecer una base sólida para el análisis y diseño del sistema, se llevó a cabo una reunión inicial de toma de requisitos el día 03/03/2026.

La entrevista se realizó con el responsable del proyecto por parte de la empresa cliente y tuvo como finalidad principal comprender el contexto del negocio, los procesos actuales de trabajo y las necesidades funcionales y no funcionales que deberá cubrir la aplicación web.

Para la sesión se utilizó una técnica de entrevista semiestructurada, apoyada en un guion previamente elaborado que permitió abordar, entre otros, los siguientes aspectos:

\begin{itemize}
    \item Objetivos del negocio y problemática actual.
    \item Flujo de trabajo existente y herramientas utilizadas.
    \item Tipos de usuarios que interactuarán con el sistema.
    \item Funcionalidades deseadas y prioridades.
    \item Restricciones técnicas y organizativas.
\end{itemize}

Con el fin de garantizar la fidelidad de la información recogida, la entrevista fue grabada en audio mediante un dispositivo Plaud Note \citep{plaudnote}, que permitió generar una transcripción automática posterior. Dicha transcripción fue revisada manualmente para asegurar la correcta interpretación de los requisitos expresados.

La información obtenida durante esta sesión ha servido como base para la elaboración de los modelos conceptuales, la definición de casos de uso y la especificación de los requisitos funcionales y no funcionales que se detallan en los apartados siguientes.